\chapter{Discussion and Limitations}

This chapter interprets the results through the same bounded-claim discipline used in evaluation. It first states what the prototype establishes, then examines technical, legal, human, and operational validity. The discussion does not repeat the result tables; it explains why each remaining limitation changes the authority of the conclusion and which future evidence would be capable of resolving it.

\section{Interpretation of Findings}

\subsection{What the Prototype Establishes}

The prototype shows that a permission-review product can be constructed so that it does not need complete visibility or automated legality in order to make a bounded statement. The core value is not the threshold alone; it is the evidence contract surrounding the threshold. Source, window, pack, rule reference, and missing context make the decision traceable. Fail-closed validation prevents malformed evidence from becoming a reassuring result. The interface then communicates the distinction between an observed signal, a synthetic example, and an evidence gap. Whether that design is in fact useful, or trusted appropriately, to any population remains an empirical question that this thesis does not answer; what is established here is the implementability and internal coherence of the design, not a measured user benefit.

The replacement of a GDPR-specific engine with regulation packs also improves scientific clarity. It reveals which statements belong to generic software and which depend on a jurisdictional interpretation. This supports future extension while making clear that interface compatibility is not legal equivalence.

\subsection{Construct Validity}

The deterministic confusion matrix measures agreement with a defined threshold oracle rather than whether that threshold identifies unlawful processing, and the other evidence types are likewise distinct analytical objects: manifest declarations, DEX references, static paths, runtime traces and policy prompts are each narrowed to what they represent.

\subsection{Internal Validity}

Test fixtures and implementation share parts of the same conceptual specification, so correlated defects remain possible. Boundary suites and the curated 9/9 mutant detections show sensitivity to nine registered weakenings but do not enumerate all plausible defects. The independent mapping packet has been prepared but not reviewed; independent implementation review, completed blinded mapping review, broader mutation generation, and event-level reference traces would further strengthen internal validity.

The simulator is intentionally useful for reproducible demonstration, but it can inflate perceived confidence if its labels are mistaken for field ground truth. The evaluated prototype mitigates this by preserving synthetic provenance through the repository and interface. The remaining risk is user over-reliance on any precise metric; future builds should consider hiding precision and recall outside an explicitly labelled research panel.

\subsection{Device Validity}

Physical evaluation used one OPPO PERM00 device and a short acceptance window. OEM process management, Android versions, accessibility settings, screen sizes, font scaling, languages, network states, and long-run scheduling can change behaviour. No claim is made for 24-hour WorkManager reliability, battery impact, memory leak freedom, or multi-device stability. A release-candidate status at this layer means the tested artefact is coherent enough for controlled store submission work, not that broad compatibility is established.

The ADB coordinate calibration problem also illustrates a methodological limitation: automated interaction evidence depends on the relationship between physical display coordinates, application-window bounds, navigation bars, and screenshots. The 24 August record anchors application actions to captured UI nodes, but the attempted extra-large-font setting could not be committed on ColorOS and no screen reader was enabled. Future UI automation should retain resource identifiers, window geometry, font state, and enabled accessibility services with every run.

\subsection{External Validity}

The \FDroidCensusN{}-APK census provides heterogeneous, hash-pinned open-source artifacts, while the \DevicePackagesN{}-package runner retains its complete eligible denominator and failures. F-Droid remains an open-source repository with its own publishing and maintenance selection effects, and the device subset is constrained by the frozen execution receipts. Neither sample supports prevalence claims for all Android applications, commercial stores, countries, users, or alleged infringements.

\subsection{Ecological Validity}

Static capability and potential source-to-sink paths do not establish executed behaviour in ordinary use. Conversely, a bounded launch workflow samples only the states reached during that run. The current evaluation therefore preserves a distinction among packaged capability, potential static flow, observed device workflow, and policy review prompt; field behaviour would require authorised event-level observation under a separately approved protocol.

\subsection{Legal and Regional Validity}

The GDPR pack maps technical signals to questions motivated by official provisions \cite{gdpr2016}. It cannot determine controller identity, processing purpose, lawful basis, Article 9 condition, necessity, proportionality, consent validity, data-subject expectations, retention, recipients, safeguards, or jurisdiction. The application's output is therefore an accountability aid. A qualified reviewer must examine the underlying processing and applicable law.

Adding another regional pack requires more than translation. The maintainer must identify scope, legal authority, effective date, regulated actors, exceptions, remedies, and concepts that have no GDPR analogue. Each pack needs independent legal sign-off, versioned fixtures, source preservation, localisation review, and an update policy. The current Research Baseline pack is not evidence that a second legal regime has been validated.

\subsection{Human Validity}

The redesigned interface applies recognised accessibility and usability principles \cite{wcag22,androidAccessibility,iso924111}, and the rendered evidence supports a bounded engineering assessment of labels, touch targets, evidence caveats, and the pause between signal and action. However, conformance-oriented inspection cannot establish comprehension, efficiency, satisfaction, trust calibration, or willingness to continue using the product, and Chapter 5 records which accessibility evidence is still missing. Those gaps are required to be closed before making an accessibility-conformance or user-adoption claim.


\subsection{Interpreting the Engineering Result}

The evaluation is strongest when read as a chain of bounded claims rather than as one pass or fail judgement. At the innermost layer, unit and integration tests show that a specified input produces the specified state transition. At the next layer, release assembly and manifest inspection show that the tested TypeScript is incorporated into an installable Android artefact with the recorded package properties. Physical-device interaction then shows that selected routes render and selected controls respond on one hardware and operating-system combination. No layer automatically establishes the next: a correct function may be absent from a package, an installable package may render poorly, and a coherent interface may still be misunderstood by users.

This layered reading prevents a common reporting error in prototype research: promoting evidence because several individually positive checks appear together. Ten successful checks are not equivalent to one check of a broader proposition. A passing threshold test and a successful device launch, for example, do not jointly prove that the threshold represents unlawful processing; they establish two different facts, one logical and one operational. Privacy Lens therefore records evidence by claim type. The report may look more cautious than a conventional demonstration, but it is more useful to a maintainer because it shows which additional study could change a conclusion.

The phrase ``initial store-submission engineering candidate'' is deliberately narrow. It means that the repository can produce the expected release formats, that the product has a coherent user journey and disclosures, and that no known in-scope engineering defect makes controlled submission work irrational. It does not mean that the application is signed with the owner's production key, that a store has reviewed its declarations, or that the legal mappings have been approved by counsel. Separating these meanings matters for privacy software because a polished security or compliance interface can borrow authority from the domain even when its evidence is limited.

There is also a difference between absence of an observed failure and affirmative evidence of a desirable property. No crash during a short device session is evidence about that session, not positive evidence of long-term stability. A small memory range across a few restarts is a useful anomaly check, not a leak proof. No sensitive runtime permission in the merged manifest is a strong package fact, but it does not by itself prove that every future acquisition connector will remain privacy preserving. The evaluation uses explicit negative claims to preserve these distinctions.

The result should therefore be understood as a reproducible decision point. A future reviewer can rebuild the same revision, inspect the same rules and fixtures, and ask whether later evidence warrants promotion from one claim class to another. That property is more durable than a one-time declaration of readiness. It converts the thesis from a retrospective narrative into a maintainable assurance record.

\section{Technical Validity Boundaries}

The central technical limitations concern the authority of acquisition, the meaning of prototype thresholds, temporal aggregation, and the governance of rule data. These issues are grouped because each can make a deterministic output reproducible while still leaving its real-world interpretation uncertain.

\subsection{Acquisition Authority and Observation Quality}

The largest external-validity constraint is not the visual interface or the threshold arithmetic; it is the authority and completeness of the observation source. Android deliberately limits one ordinary application's access to sensitive information about other applications. Privacy Lens does not bypass that security model. The current native bridge and imported-record contracts can carry observations when an authorised deployment supplies them, while the simulator supports deterministic explanation and testing. The prototype must not imply that installing it from a public store grants unrestricted historical visibility.

This constraint affects every downstream conclusion. An apparently quiet history may mean that accesses did not occur, that the selected API does not expose them, that an OEM shortened retention, that the collection job was delayed, or that the current deployment lacks authority. These explanations are not interchangeable. The evidence schema therefore includes source and time-window fields, and an evidence-gap state is preferred when acquisition sufficiency is unknown. A ``no technical concern'' result refers only to accepted records, not to all activity on the device.

Observation quality has at least five dimensions: coverage (which applications, operations and intervals are visible), resolution (event-level or aggregated), timeliness, integrity (whether records can be altered, replayed or reordered), and semantic fidelity (whether the operation label corresponds to the personal-data behaviour under review). A future connector should publish a capability document addressing all five, because merely conforming to the TypeScript shape is insufficient. Aggregate records trade context for storage — ten accesses could be ten user-requested refreshes or repeated background tracking, and a total cannot distinguish them — so the acquisition design must minimise collected fields, bound retention and protect local storage. More data is not automatically a more privacy-preserving audit.

An authorised enterprise or research deployment could improve coverage through device-owner controls, instrumented test applications or controlled firmware, but each option changes the population to which the findings apply, so evaluation reports must name the authority model rather than treat all Android evidence as equivalent. The validation programme that would turn coverage into an empirical statement is specified in Chapter 7.

\subsection{Threshold Meaning, Calibration, and Decision Error}

The implemented baselines and multipliers are transparent prototype parameters, as Chapter 3 records. Transparency makes a rule reproducible without making it representative: a location threshold can be inspected and tested exactly, but its normative meaning depends on application category, user action, service duration and declared purpose, so a navigation application and a static calculator cannot be judged by the same behavioural expectation merely because both run on Android.

False positives and false negatives also have asymmetric consequences. A false positive may alarm a user, damage confidence in a legitimate application, and contribute to warning fatigue. A false negative may reassure the user or reviewer despite activity that deserves investigation. The appropriate balance is not a purely mathematical optimisation because the costs depend on use context. The current design therefore avoids converting a below-threshold result into a compliance statement and avoids converting an exceedance into a violation verdict.

The calibration target must be explicit, because optimising agreement with human reviewers is not the same as predicting store-policy enforcement, detecting malware, or identifying GDPR accountability questions; each target requires different labels and evidence. Privacy Lens targets review prioritisation — a defensible prompt for further examination — which supports conservative wording and structured missing-context questions rather than a binary compliance classifier. The sampling frame, annotation protocol, multi-scale parameter governance and reporting stratification that a calibration programme would need are specified in Chapter 7.

\subsection{Temporal Evidence and Replay Semantics}

The failure modes of time-windowed evidence, and the controls the engine applies to them, are set out in Chapter 3. Two consequences belong in the discussion rather than in the design chapter. The controls specify how the engine behaves under replay without creating event-level identity where none exists, so partially overlapping aggregates cannot always be reconciled; and the conservative response — declining cumulative interpretation when independence cannot be established — trades sensitivity for the avoidance of manufactured precision. Retention follows the same reasoning: keeping only the most recent records bounds memory and reduces local exposure, yet an arbitrary truncation point can remove context needed for a decision, so the history limit is an engineering safeguard rather than a legal retention determination. The temporal properties the registered suite should additionally cover, and the connector designs that would strengthen deduplication, are listed in Chapter 7.

The broader lesson is that time is part of the evidence contract, not metadata attached after classification. A risk statement about repeated behaviour is only as defensible as the window relationships used to compute it. Making those relationships visible helps reviewers distinguish sustained behaviour from an artefact of aggregation.

\subsection{Regulation-Pack Governance Beyond Software Modularity}

The replaceable pack contract solves a software architecture problem: it prevents jurisdiction-specific references, prompts and thresholds from being scattered across screens, storage and native integration code. It does not solve the institutional problem of deciding what a pack should say, because a pack can be syntactically valid and legally misleading. Governance therefore has to be paired with the architecture on two fronts.

The first is admission. A pack must name its jurisdiction, regulated context, legal instruments, effective date, intended users and exclusions; distinguish binding provisions from regulator guidance and from research heuristics; and link each technical condition to a review question plus the missing facts that prevent automated legal determination. Legal review must be independent of the developer who encoded the rule, must cover citation accuracy, applicability, exceptions, terminology and translations, and must attach to a specific pack version and date. Technical admission then validates identifiers, semantic versions, rule uniqueness, supported evidence fields, bounded parameters, localisation completeness and fixture coverage, and the registry rejects an unknown or malformed pack rather than falling back silently to a default. Chapter 3 records the contract itself; Chapter 7 lists the deprecation, rollback and offline-recovery machinery that a downloadable pack would additionally require.

The second is regional variation. Concepts such as controller, covered business, sensitive information, consent, sale, sharing, child or data localisation do not map one-to-one, and enforcement actors and private rights differ, so a common model must carry generic evidence and review primitives while letting packs declare jurisdiction-specific questions; forcing every regime into GDPR terminology would produce superficial flexibility and substantive error. Localisation is likewise more than translating strings: it combines legal-language review, plain-language review in the user's locale, and rendered interface testing, and the interface must identify both the selected region and the pack version. The Research Baseline pack illustrates the disclosure pattern the thesis relies on — switching rule sets while preserving identity, without masquerading as a validated regional law.

\section{Human, Ethical, and Operational Trust}

The product's credibility depends on more than a correct rule. Users must understand the status, the reviewing application must minimise its own privacy and security risks, research evidence must be collected responsibly, and release claims must match the package that is actually distributed.

\subsection{Trust, Comprehension, and Willingness to Use}

Visual coherence is necessary because users make rapid judgements about whether a privacy tool is maintained and safe. Consistent typography, restrained colour, stable navigation, and plain-language hierarchy reduce avoidable friction. They are not sufficient for trustworthiness. A visually polished product can be more dangerous if polish causes users to accept unsupported conclusions. Privacy Lens therefore treats calibrated trust, rather than maximum trust, as the design objective.

Calibrated trust means that confidence changes with evidence quality. A native or authorised imported observation can be described according to its documented capability. A simulator record must remain labelled synthetic. Missing purpose or lawful-basis context must remain visible even when the numerical signal is normal. The interface should make the most important limitation available at the decision point, not bury it exclusively in a policy page.

Willingness to use has several components. Initial adoption depends on a credible value proposition and low setup burden. Continued use depends on whether findings arrive at useful times, whether repeated alerts are suppressed, whether explanations support action, and whether local data controls behave predictably. Recommendation to others may depend on institutional endorsement, independent review, and evidence that the tool does not create new privacy risks. A short expert inspection can assess obvious barriers but cannot establish these behavioural outcomes.

A participant study should therefore use realistic tasks rather than preference ratings alone. Participants could identify the source of a finding, compare two risk states, explain why a low-count record is not a compliance certificate, switch regulation packs, clear local evidence, and decide what follow-up information is needed. Measures should include completion, error, time, comprehension, perceived workload, appropriate reliance, and confidence. Interviews can reveal whether terms such as evidence gap or technical concern match users' mental models.

The sample should include ordinary Android users, privacy-aware users, accessibility participants, and professional reviewers. Their goals differ. Ordinary users may need a concise decision path; professionals may need exportable provenance and exact rule details. Accessibility participants may expose focus order, verbosity, contrast, magnification, or motor-control barriers that visual review misses. The study should avoid treating one group's preference as universal.

Longitudinal testing is necessary for notification design. A single session cannot reproduce fatigue or habituation. A multi-week study could compare immediate alerts, daily summaries, and state-change-only notifications. It should record dismissals and follow-up actions while avoiding collection of unnecessary behavioural telemetry. The hypothesis is not that fewer notifications are always better, but that each interruption should correspond to a meaningful change.

Trust measures also need a correct-answer component. Asking whether participants trust the application can reward unjustified confidence. Better evaluation presents cases with strong evidence, synthetic evidence, missing context, and malformed input, then tests whether confidence tracks those differences. An effective design should sometimes cause the user to withhold judgement. That is a successful outcome for an evidence-bounded review tool.

\subsection{Privacy and Security of the Reviewing Application}

A privacy-review application can itself become a source of sensitive data. Permission histories may reveal health routines, mobility, communication, or working patterns even when they contain no content. Application identifiers can expose relationships or interests. Findings and notification text may be visible to other people with physical access to the device. The product's local-first design reduces network disclosure, but local processing is not automatically risk free.

The current package requests no sensitive runtime permission and disables Android backup. These are valuable minimisation facts for the evaluated build. The application also provides a clear-data control and does not require an account. Future connectors could change this profile. Every new acquisition or synchronisation feature should therefore trigger a new data-flow inventory, manifest reconciliation, threat review, privacy notice update, and store declaration review.

At-rest protection should be proportionate to the threat model. Platform sandboxing protects data from ordinary applications, while device compromise, debug builds, rooted environments, and unlocked backups present different risks. If future deployments store event-level histories or legal annotations, encryption and key lifecycle may become appropriate. Encryption does not eliminate exposure through rendered screens, notifications, logs, or exported files, so those channels also require controls.

Logging is a frequent source of accidental disclosure. Production builds should avoid writing raw evidence, application inventories, or legal-context notes to Logcat. Crash reports, if introduced, should be opt-in where required, minimise payloads, document processors and retention, and avoid attaching the evidence database. Analytics should not be added merely to measure engagement; a privacy tool carries a higher burden to justify behavioural tracking.

Pack supply-chain risk must also be considered. A malicious or compromised pack could present alarming text, suppress review, or direct users to unsafe resources even without executable code. Declarative schema restrictions, signed distribution, trusted publisher identities, URL allow-lists where appropriate, and human review reduce this risk. Bundled packs should be covered by the same source-control and release-review process as application code.

Abuse scenarios include importing fabricated evidence to defame an application, replaying old records as current, selecting a misleading jurisdiction, and showing a screenshot without its source label. Privacy Lens addresses some of these risks through provenance, timestamps, pack identity, and replay controls. It cannot authenticate every external producer in the current prototype. Export formats should preserve those fields and should state that the document is a review artefact rather than a regulator finding.

Security maintenance is part of store readiness. Dependencies, Android target requirements, signing keys, build provenance, and vulnerability disclosures require ongoing ownership after submission. The present thesis records a reproducible snapshot; it cannot guarantee that the same dependencies remain appropriate later. A release process should include dependency review, secret scanning, reproducible build records, and an incident contact before the product is distributed broadly.

\subsection{Research Ethics and Responsible Evaluation}

The safest current evaluation uses synthetic or controlled evidence because real permission histories can expose third-party and participant behaviour. Synthetic data provides precise ground truth and repeatability, but it must be labelled so screenshots and metrics are not mistaken for field evidence. The simulator label is therefore an ethical as well as methodological control.

A future participant or device study would require informed consent describing what is collected, why it is needed, how long it is retained, who can access it, and how withdrawal operates. Researchers should avoid collecting unrelated application inventories or fine-grained histories when controlled test applications can answer the research question. Test accounts and scripted actions should be preferred for validation of acquisition correctness.

The legal subject matter creates an additional risk of participant misunderstanding. Study materials should explain that findings are prompts and that the research team is not providing individual legal advice. Participants should not be encouraged to confront application developers or change essential device settings solely because of a synthetic warning. Debriefing should clarify which cases were fabricated and how the evidence boundaries work.

Reporting should include failed runs, missing observations, exclusions, and changes made after pilot testing. Selectively presenting only coherent screenshots would exaggerate maturity. The repository's separation of logical, package, physical, and visual evidence provides a structure for ethical reporting because each claim can be tied to its acquisition method.

If external legal experts annotate cases, their independence, jurisdiction, instructions, and disagreement process should be recorded. A single expert opinion is valuable but should not be transformed into universal ground truth. Where reasonable interpretations differ, the pack can present the ambiguity as a question rather than forcing a binary label.

Finally, public release creates responsibilities to people who were not research participants. Store text and onboarding must not imply regulatory endorsement. Support channels should handle correction and deletion questions. Pack updates should communicate material interpretation changes. Responsible deployment is therefore a continuing governance activity, not the final administrative step after technical acceptance.

\subsection{Store and Operational Readiness}

The repository contains target-SDK configuration, AAB output, privacy-policy content, a stable package identifier, branded assets, and a store checklist. Google Play nevertheless requires owner-controlled signing and console information, including a Data safety declaration and current policy compliance \cite{playDataSafety,playTargetApiLevel}. The present AAB uses a QA signing fallback when no upload-key environment is provided. Public submission therefore remains blocked on an upload key, HTTPS privacy-policy publication, support contact, final listing content, screenshots selected for the listing, console declarations, internal-track testing, and store review.

\section{Temporal-Alert Limits and Next Validation}

Temporal co-occurrence is necessarily an incomplete proxy. A combined location, sensor, or transfer pattern may be benign, expected by the user, technically misclassified, partially observed, or governed by a lawful purpose not present in local evidence. Conversely, an actual high-risk flow can evade this prototype when it occurs outside a profile's window, uses an unsupported interface, is hidden by platform restrictions, or lacks observable bridge evidence. The bounded ledger intentionally trades historical completeness for local resource control. These are reasons to preserve uncertainty and notification-only operation, not reasons to silently broaden a rule.

Future packs should be authored through controlled legal and engineering review, with versioned source material, traceable interpretation notes, fixtures for matches and non-matches, and a staged rollback path; the rule author rather than the generic engine must set and justify each window and required combination, and pack loading must keep validating the converter boundary so that unsupported observation types, empty legal references, unsafe counts or implausible windows cannot activate accidental alerts. The external validation that any broader claim would need is listed in Chapter 7.

\section{Residual Limits}

Persistence improves continuity only within one installed application and the active maximum window. It is not tamper resistance, backup continuity, a chain of custody, or a substitute for an authorised observation source. Clock rollback, data clearing, uninstallation, damaged storage, pack changes, and operating-system limits remain fail-closed boundaries. The controlled route now has one OPPO PERM00 execution record verifying debug-native transport, temporal evaluation, ledger update, and presentation; it does not verify sensing, collection, third-party AppOps history, a legal violation, or any other device. CI demonstrates repeatable engineering regression checks; it does not replace device diversity, independent legal assessment, or field validation.

\section{Prioritised Next Work}

The next programme should proceed in this order:

\begin{enumerate}
    \item obtain independent GDPR-pack review and define pack-update governance;
    \item provision production signing and complete Play Console privacy and Data safety declarations;
    \item run accessibility and localisation matrices, followed by task-based participant evaluation;
    \item execute multi-OEM, multi-version, 24-hour scheduling, battery, memory, and recovery tests;
    \item replace aggregate-only acquisition where authorised event-level evidence is available;
    \item design and independently validate one additional regional pack before claiming multi-jurisdiction support.
\end{enumerate}

This order prioritises validity and operational trust over adding more visible features.

